\section{Auxiliary results}\label{sec:aux}

The development contains a number of results that were built as tools for the main theorem or as
tests of the frontier and stand on their own. This section lists the ones an economist may want.

\paragraph{Multiplicity is exactly non-monotone supply, and the equilibrium system permits it.}
Two distinct equilibrium rates force capital supply to be non-monotone on any interval containing
them (\lean{not\_monotoneOn\_of\_equilibria\_ne}), so every multiplicity example in the literature
is a counterexample to Light's Theorem 1 and to nothing else. Conversely an explicit continuous,
strictly positive supply schedule, Cobb--Douglas demand plus a cubic wiggle, has exactly three
equilibria with the demand $D(r)=A^2/(4(r+\delta)^2)$ (\lean{wiggleSupply\_equilibria}),
matching the three-equilibrium shape \citet{YoungWalsh2025} report: nothing in the firm side, the
fixed-point machinery or the intermediate-value argument rules multiplicity out. What is not
proved is that such a schedule is the capital supply of any household; the two bounds available,
the gain test and the minimal MPC, both move the wrong way for a backward bend.

\paragraph{All capital is precautionary.}
A two-state economy near log utility holds at least $1/400$ of capital at every stationary
distribution. The same economy with its income risk removed, both states replaced by their mean,
saves nothing at any asset level, so its capital supply is exactly zero
(\lean{riskless\_aggregateCapital\_eq\_zero}). The corner condition of Section~\ref{sec:corner}
holds at every asset level in the riskless economy and fails above a threshold in the risky one,
and the gap is the value of the insurance the household is buying.

\paragraph{Impatience is necessary.}
\citet{MarcetObiolsWeil2007}, Proposition 3: no stationary equilibrium with finite capital exists at
$\beta R\ge1$. Formalised with income risk through the Euler route, requiring $\beta R>1$ strictly,
and deterministically through secant slopes (\lean{Impatience}); the version with a supply floor is
Theorem~\ref{thm:above}.

\paragraph{Carroll--Kimball concavity across HARA.}
The consumption function is concave in assets for CRRA, shifted CRRA (via a power-mean
inequality), and CARA (via a soft-minimum construction); quadratic utility is proved inadmissible
because it is not increasing (\lean{not\_monotoneOn\_quadUtility}). The induction runs along
Bellman iterates, so it needs the Euler equation one step at a time and at states where the
constraint binds, which is what \lean{IncomeFluctuationEuler} supplies. A witness shows Carroll and
Kimball's conclusion is false when the cap binds (\lean{not\_concaveOn\_consumptionFn\_of\_cap\_binds}),
which is why the cap slack condition cannot be dropped.

\paragraph{A calibrated class with existence and uniqueness.}
Before the log calibration was reached, existence and uniqueness were proved for a class of
economies (\lean{Calibrated}) each of whose numerical fields is one primitive inequality: near-log
CRRA, log Stone--Geary (\lean{stoneGeary\_exists\_equilibrium\_of\_technology'}), and the $b\neq0$
branch of HARA, with general finite Markov income. These carry the machinery of
Section~\ref{sec:arch} at parameters where the published constants do work, and a CES witness
proved uncalibratable shows the cap conditions are not vacuous.

\paragraph{Sufficiency of the Euler equation, without differentiability of the value function.}
A solution of the functional Euler equation with equality at interior choices and the correct
inequality at the constraint is the optimal consumption function
(\lean{eq\_consumptionFn\_of\_isEulerSolution}). The proof is the contraction of the Coleman
operator in the marginal-utility supremum metric of \citet{LiStachurski2014}, run on a class of
candidate policies with a consumption floor and a saving slack. This is the sufficient direction
of the Coleman--Reffett conjugacy of \citet{SargentStachurski2024}, Proposition 8.3.13, proved
without the conjugacy.

\paragraph{Abstract fixed-point results from Sargent and Stachurski.}
Parametric monotonicity of fixed points (Vol.\ 1, Prop.\ 2.2.7); eventually contracting maps
(Thm.\ 6.1.5), which is what spectral-radius conditions deliver; the fitted-value-iteration error
bound (Vol.\ 2, Thm.\ 9.1.3), the value-error input to the policy error bounds of \citet{Li2015};
and uniqueness of fixed points of increasing concave maps, in one dimension and in Du's form on
a finite product (\lean{OrderedFixedPoints}). \citet{Santos2000}'s policy error bound was
assessed and not formalised: it needs interior solutions, which the borrowing constraint denies.

\paragraph{Mean-field-game uniqueness, and why it does not reach Bewley.}
The Lasry--Lions argument for uniqueness of mean-field equilibria adds two optimality inequalities.
Its measure-theoretic core is formalised for the static game and for the finite-horizon
finite-state game on paths: two mean-field Nash distributions have nonpositive Lasry--Lions pairing;
under monotone costs the pairing vanishes; under strict monotonicity the costs coincide, though the
distributions need not (\lean{cost\_eq\_of\_nash}, \lean{value\_eq\_of\_pathEquilibrium}; the
static form of \citealp{CannarsaCapuani2018}, Thm.\ 4.1). The argument cancels
$\int(u_1-u_2)\,d(m_1(0)-m_2(0))$ by a common initial distribution, which a stationary discounted
equilibrium does not have; and \citet{GraberMatter2024} show that games coupled through a market
price are typically not Lasry--Lions monotone, the right condition there being monotonicity of the
error map $Q\mapsto Q-\text{best response}(Q)$, which for Aiyagari is Light's single crossing. The
quantitative version, a strongly monotone Lipschitz operator on a real inner product space has a
unique zero found by a damped iteration with modulus $\sqrt{1-m^2/L^2}$, is formalised
(\lean{exists\_unique\_zero\_of\_stronglyMonotone}); for excess supply it would need a slope floor
and a Lipschitz constant in the rate, neither of which is available.

\paragraph{Tauchen's method.}
The Tauchen transition matrix is defined from cumulative rows of the standard normal distribution
function, taken from Mathlib's Gaussian measure, and proved to be stochastic, positive in its first
column, and monotone in the sense of \ref{E1} for $\rho\ge0$
(\lean{monotoneTransitions\_of\_tauchen}). This is what allows Corollary~\ref{cor:tauchen} to be
stated for any number of grid points.
